\section{Thinking and declaration}
This project highlights the broad potential of AI-driven MRI reconstruction in accelerating imaging, improving reconstruction quality, and reducing acquisition cost. By combining deep learning with multimodal guidance and frequency-domain modeling, modern reconstruction systems can better exploit both anatomical priors and physical measurement information. Our experiments suggest that future MRI reconstruction methods may continue to improve in efficiency, image fidelity, and clinical applicability, especially when image-domain learning is more tightly integrated with k-space supervision and modality-aware fusion strategies.

During the development of this project, we used AI-assisted tools and large language models, including Codex and Gemini, to support implementation and coding-related tasks. We provided the core ideas, methodological design, experimental goals, and technical requirements, while the AI systems assisted with code generation, implementation suggestions, and debugging support. All AI-generated outputs were carefully reviewed, validated, and iteratively refined by us before being incorporated into the final project. We further modified and integrated these outputs to ensure their correctness, reliability, and consistency with the objectives and academic standards of this work.

\section{Acknowledgement}
We would like to express our sincere gratitude to Professor Qi for the course guidance, insightful suggestions, and continuous support throughout this project. His instruction helped us better understand both the medical imaging background and the methodological challenges of MRI reconstruction. We also sincerely thank the teaching assistants for their patience, helpful feedback, and timely support during the implementation and experimentation process. Their assistance was very valuable to the completion of this work.
